\chapter{Schémas de Lewis, mésomérie et VSEPR}\label{ch:b1:lewis-resonance-vsepr}

La molécule d'ozone, \ce{O3}, est une chaîne coudée de trois atomes
d'oxygène. Dessinée avec les règles du livre 1, elle présente une liaison
double et une liaison simple~: l'une devrait être courte, l'autre longue.
Pourtant, la spectroscopie micro-ondes trouve les deux liaisons \ce{O-O}
de l'ozone rigoureusement égales, à \qty{127.8}{pm}, entre la liaison
double de \ce{O2} (\qty{120.8}{pm}) et la liaison simple du peroxyde
d'hydrogène (\qty{147.5}{pm}). Aucun dessin isolé n'est juste~; deux
dessins ensemble le sont. Ce chapitre précise le modèle de Lewis ---
\omterm{def:b1:lewis-resonance-vsepr:formal-charge}{charges formelles}, octets incomplets ou dépassés, mésomérie entre
plusieurs formules --- puis s'en sert pour prévoir la géométrie des
molécules et les \omterm{def:b1:lewis-resonance-vsepr:dipole}{moments dipolaires} qui en découlent.
% ledger: geo:O3.rOO, re:O2, geo:H2O2.rOO

\begin{recall}
Le livre 1 (année 10) a décrit la \omterm{def:b1:lewis-resonance-vsepr:lewis}{liaison covalente} comme un doublet
d'électrons mis en commun, a dessiné des \omterm{def:b1:lewis-resonance-vsepr:lewis}{schémas de Lewis} avec \omterm{def:b1:lewis-resonance-vsepr:lewis}{doublets
liants} et non liants, et a prévu quatre géométries (linéaire, coudée,
triangulaire, tétraédrique). L'\omterm{def:b1:periodicity:electronegativity}{électronégativité} et ses différences ont
été rendues quantitatives au \cref{ch:b1:periodicity}~; les \omterm{def:b1:quantum-numbers:configuration}{électrons de
valence} d'un atome se lisent sur sa configuration
(\cref{def:b1:quantum-numbers:configuration}).
\end{recall}

\section{Schémas de Lewis et charges formelles}

\begin{definition}[Liaison covalente, schéma de Lewis]\label{def:b1:lewis-resonance-vsepr:lewis}
Une \emph{liaison covalente}\index{liaison covalente} est un doublet d'électrons mis en
commun par deux atomes, le \emph{doublet liant}\index{doublet liant}~; deux ou trois doublets
forment une liaison double ou triple. Un doublet d'\omterm{def:b1:quantum-numbers:configuration}{électrons de valence}
qui appartient à un seul atome est un \emph{doublet non liant}\index{doublet non liant}. Un
\emph{schéma de Lewis}\index{schema de Lewis@schéma de Lewis} représente tous les \omterm{def:b1:quantum-numbers:configuration}{électrons de valence}
d'une molécule ou d'un ion, sous forme de liaisons (traits entre les
atomes) et de doublets non liants (traits ou paires de points sur un
atome). Selon la \emph{règle de
l'octet}\index{regle de l'octet@règle de l'octet}, les atomes de la \omterm{def:b1:periodicity:table}{période} 2 (C, N, O, F) sont
entourés, dans une structure stable, de quatre doublets, soit huit
électrons~; selon la \emph{règle du duet}\index{regle du duet@règle du duet}, l'hydrogène l'est d'un
seul doublet.
\end{definition}

\begin{definition}[Charge formelle]\label{def:b1:lewis-resonance-vsepr:formal-charge}
La \emph{charge formelle}\index{charge formelle} d'un atome dans un \omterm{def:b1:lewis-resonance-vsepr:lewis}{schéma de Lewis}
vaut
\[
  q_F = N_v - N_{\text{nl}} - \tfrac12 N_{\text{l}},
\]
où $N_v$ est le nombre d'\omterm{def:b1:quantum-numbers:configuration}{électrons de valence} de l'atome libre,
$N_{\text{nl}}$ le nombre d'électrons de ses \omterm{def:b1:lewis-resonance-vsepr:lewis}{doublets non liants} et
$N_{\text{l}}$ le nombre d'électrons des liaisons qu'il forme. On
l'écrit $\oplus$ ou $\ominus$ à côté de l'atome lorsqu'elle n'est pas nulle.
\end{definition}

\begin{proposition}[La somme des charges formelles est la charge]\label{prop:b1:lewis-resonance-vsepr:formal-sum}
La somme des \omterm{def:b1:lewis-resonance-vsepr:formal-charge}{charges formelles} des atomes d'un \omterm{def:b1:lewis-resonance-vsepr:lewis}{schéma de Lewis} est égale
à la charge totale de la molécule ou de l'ion.
\end{proposition}

\begin{proof}
La somme de $q_F$ sur les atomes vaut $\sum N_v - (\text{électrons des
doublets non liants} + \text{électrons des liaisons})$, car les électrons
de chaque liaison sont comptés pour moitié sur chacun de ses deux atomes.
La parenthèse est le nombre total d'\omterm{def:b1:quantum-numbers:configuration}{électrons de valence} dessinés dans le
schéma, égal à $\sum N_v$ moins la charge $z$ de l'espèce. Donc
$\sum q_F = z$.
\end{proof}

\begin{method}[Dessiner un schéma de Lewis]\label{met:b1:lewis-resonance-vsepr:lewis}
Pour dessiner le \omterm{def:b1:lewis-resonance-vsepr:lewis}{schéma de Lewis} d'une molécule ou d'un ion~:
\begin{enumerate}
  \item compter les \omterm{def:b1:quantum-numbers:configuration}{électrons de valence} $N = \sum N_v - z$ et les
        doublets $N/2$~;
  \item dessiner le squelette~: l'atome le moins électronégatif (autre que
        H) est en général central~; l'hydrogène et le fluor sont toujours
        terminaux~;
  \item compléter les octets des atomes terminaux par des \omterm{def:b1:lewis-resonance-vsepr:lewis}{doublets non
        liants}, puis placer les doublets restants sur l'atome central~;
  \item si l'atome central n'a pas son octet, transformer des \omterm{def:b1:lewis-resonance-vsepr:lewis}{doublets
        non liants} de ses voisins en liaisons multiples~;
  \item calculer les \omterm{def:b1:lewis-resonance-vsepr:formal-charge}{charges formelles}~; parmi les structures possibles,
        préférer celle qui porte le moins de \omterm{def:b1:lewis-resonance-vsepr:formal-charge}{charges formelles}, avec les
        charges négatives sur les atomes les plus électronégatifs.
\end{enumerate}
\end{method}

\begin{example}[L'acide nitrique]\label{ex:b1:lewis-resonance-vsepr:hno3}
\ce{HNO3}~: $N = 1 + 5 + 3 \times 6 = 24$ électrons, 12 doublets.
Squelette~: l'azote central, lié à trois oxygènes dont l'un porte
l'hydrogène. Compléter les octets avec des liaisons simples laisse
l'azote avec six électrons~; un \omterm{def:b1:lewis-resonance-vsepr:lewis}{doublet non liant} d'un oxygène devient
une liaison \ce{N=O}. \omterm{def:b1:lewis-resonance-vsepr:formal-charge}{Charges formelles}~: l'azote $5 - 0 - 4 = +1$~;
l'oxygène lié par une liaison simple et sans hydrogène $6 - 6 - 1 = -1$~;
les autres 0. Le total est nul, comme l'exige la
\cref{prop:b1:lewis-resonance-vsepr:formal-sum}.
\end{example}

\begin{omfigure}
\setchemfig{atom sep=2.4em}
\begin{tikzpicture}[font=\small]
  \node at (0,0) {\large\chemfig{H-\charge{90=\|,270=\|}{O}-\charge{90:2pt=$\scriptstyle\oplus$}{N}(=[:45]\charge{0=\|,90=\|}{O})-[:-45]\charge{0=\|,240=\|,300=\|,45:3pt=$\scriptstyle\ominus$}{O}}};
  \node[below] at (0,-1.5) {acide nitrique, \ce{HNO3}};
  \node at (5.2,0) {\large\chemfig{\charge{180=\|}{N}~\charge{0=\|}{N}}};
  \node[below] at (5.2,-1.5) {diazote, \ce{N2}};
  \node at (9.6,0) {\large\chemfig{\charge{135=\|,225=\|,90:2pt=$\scriptstyle\ominus$}{N}=\charge{90:2pt=$\scriptstyle\oplus$}{N}=\charge{45=\|,315=\|}{O}}};
  \node[below] at (9.6,-1.5) {monoxyde de diazote, \ce{N2O}};
\end{tikzpicture}

{\small \omterm{def:b1:lewis-resonance-vsepr:lewis}{Schémas de Lewis} avec \omterm{def:b1:lewis-resonance-vsepr:lewis}{doublets non liants} et \omterm{def:b1:lewis-resonance-vsepr:formal-charge}{charges formelles}.}
\end{omfigure}

\section{Au-delà de l'octet}

\begin{definition}[Molécules déficitaires en électrons et hypervalentes]\label{def:b1:lewis-resonance-vsepr:hypervalent}
Une molécule dans laquelle un atome a moins de huit \omterm{def:b1:quantum-numbers:configuration}{électrons de valence}
est \emph{déficitaire en électrons}\index{deficitaire en electrons@déficitaire en électrons}~: dans \ce{BF3}, le bore en a
six. Une molécule dans laquelle un atome de la \omterm{def:b1:periodicity:table}{période} 3 ou au-delà est
entouré de plus de quatre doublets est \emph{hypervalente}\index{hypervalent}~: dans
\ce{PCl5}, le phosphore forme cinq liaisons, dans \ce{SF6} le soufre six.
Les molécules à nombre impair d'électrons, comme \ce{NO} et \ce{NO2}, ne
peuvent satisfaire la \omterm{def:b1:lewis-resonance-vsepr:lewis}{règle de l'octet} sur chaque atome~: un électron est
célibataire.
\end{definition}

\begin{remark}[Pourquoi la période 3 et non la période 2]\label{rem:b1:lewis-resonance-vsepr:hypervalent}
L'azote forme \ce{NF3} mais jamais \ce{NF5}, alors que le phosphore forme
\ce{PF3} comme \ce{PF5}. Un atome de la \omterm{def:b1:periodicity:table}{période} 2 est trop petit pour se
lier à plus de quatre voisins~; un atome de la \omterm{def:b1:periodicity:table}{période} 3 est assez gros
pour en accueillir cinq ou six. On trouve pour \ce{SO4^{2-}} ou
\ce{H2SO4} des dessins avec des liaisons doubles sur le soufre (12
électrons autour de S) et d'autres avec des liaisons simples et des
\omterm{def:b1:lewis-resonance-vsepr:formal-charge}{charges formelles} (un octet)~; le second rend mieux compte de la
répartition des charges, et tous deux donnent la bonne géométrie.
\end{remark}

\begin{definition}[Acide de Lewis, base de Lewis, liaison dative]\label{def:b1:lewis-resonance-vsepr:lewis-acid}
Un \emph{acide de Lewis}\index{acide de Lewis} est une espèce capable d'accepter un doublet
d'électrons dans une place vacante de sa couche de valence (\ce{BF3},
\ce{AlCl3}, \ce{H+}, cations métalliques)~; une \emph{base de Lewis}\index{base de Lewis}
est une espèce qui possède un \omterm{def:b1:lewis-resonance-vsepr:lewis}{doublet non liant} qu'elle peut mettre en
commun (\ce{NH3}, \ce{H2O}, \ce{F-}). Quand la base cède son doublet à
l'acide, la liaison formée est une
\emph{liaison dative}\index{liaison dative}~: une \omterm{def:b1:lewis-resonance-vsepr:lewis}{liaison covalente} dont les deux électrons
proviennent tous deux d'un même partenaire.
\end{definition}

\begin{example}[L'ammoniac et le trifluorure de bore]\label{ex:b1:lewis-resonance-vsepr:adduct}
\ce{NH3 + BF3 -> H3NBF3}. Dans l'adduit, le bore complète son octet~;
l'azote, qui met désormais en commun son \omterm{def:b1:lewis-resonance-vsepr:lewis}{doublet non liant}, porte la
\omterm{def:b1:lewis-resonance-vsepr:formal-charge}{charge formelle} $+1$ et le bore $-1$. Une fois formée, la liaison
\ce{N-B} est une \omterm{def:b1:lewis-resonance-vsepr:lewis}{liaison covalente} ordinaire.
\end{example}

\begin{omfigure}
\setchemfig{atom sep=2.1em}
\begin{tikzpicture}[font=\small]
  \node at (0,0) {\large\chemfig{H-\charge{90=\|}{N}(-[:240]H)(-[:300]H)}};
  \node at (1.6,0) {$+$};
  \node at (3.1,0) {\large\chemfig{B(-[:90]F)(-[:210]F)(-[:330]F)}};
  \draw[-{Stealth}, thick] (4.3,0) -- (5.3,0);
  \node at (7.4,0) {\large\chemfig{H-\charge{45:1pt=$\scriptstyle\oplus$}{N}(-[:90]H)(-[:270]H)-\charge{45:1pt=$\scriptstyle\ominus$}{B}(-[:90]F)(-[:270]F)-F}};
\end{tikzpicture}

{\small Une \omterm{def:b1:lewis-resonance-vsepr:lewis-acid}{liaison dative}~: le \omterm{def:b1:lewis-resonance-vsepr:lewis}{doublet non liant} de l'ammoniac (une \omterm{def:b1:lewis-resonance-vsepr:lewis-acid}{base
de Lewis}) occupe la place vacante du trifluorure de bore (un \omterm{def:b1:lewis-resonance-vsepr:lewis-acid}{acide de
Lewis}). Les \omterm{def:b1:lewis-resonance-vsepr:lewis}{doublets non liants} du fluor ne sont pas dessinés.}
\end{omfigure}

\section{Mésomérie}

\begin{definition}[Formules mésomères, hybride de résonance]\label{def:b1:lewis-resonance-vsepr:resonance}
Lorsque les électrons d'une molécule peuvent être placés selon plusieurs
\omterm{def:b1:lewis-resonance-vsepr:lewis}{schémas de Lewis} qui ne diffèrent que par la position des liaisons $\pi$
et des \omterm{def:b1:lewis-resonance-vsepr:lewis}{doublets non liants} --- les noyaux restant en place ---, chacun
est une \emph{formule
mésomère}\index{formule mesomere@formule mésomère} (ou forme limite de résonance), reliée aux autres
par une flèche à deux pointes $\leftrightarrow$. La molécule réelle n'est
aucune d'elles~: c'est l'\emph{hybride de résonance}\index{hybride de resonance@hybride de résonance}, une structure
unique dont la répartition électronique est une moyenne pondérée des
leurs.
\end{definition}

\begin{remark}[Une flèche, pas un équilibre]\label{rem:b1:lewis-resonance-vsepr:arrow}
La flèche $\leftrightarrow$ ne décrit pas une réaction~: l'ozone ne
bascule pas d'une formule à l'autre. L'hybride est une seule molécule,
décrite par plusieurs dessins parce qu'un seul dessin ne suffit pas. Ne
jamais utiliser $\rightleftharpoons$ entre des \omterm{def:b1:lewis-resonance-vsepr:resonance}{formules mésomères}.
\end{remark}

\begin{method}[Écrire et pondérer des formules mésomères]\label{met:b1:lewis-resonance-vsepr:resonance}
Pour trouver les \omterm{def:b1:lewis-resonance-vsepr:resonance}{formules mésomères} d'une espèce~:
\begin{enumerate}
  \item partir d'un \omterm{def:b1:lewis-resonance-vsepr:lewis}{schéma de Lewis}~; chercher un \omterm{def:b1:lewis-resonance-vsepr:lewis}{doublet non liant} ou
        une liaison $\pi$ voisins d'une liaison $\pi$, d'une place vacante
        ou d'une charge positive~;
  \item déplacer des doublets avec des flèches courbes (un \omterm{def:b1:lewis-resonance-vsepr:lewis}{doublet non
        liant} devient une liaison $\pi$, une liaison $\pi$ devient un
        \omterm{def:b1:lewis-resonance-vsepr:lewis}{doublet non liant}), sans jamais déplacer un noyau ni dépasser
        l'octet d'un atome de la \omterm{def:b1:periodicity:table}{période} 2~;
  \item recalculer les \omterm{def:b1:lewis-resonance-vsepr:formal-charge}{charges formelles}~;
  \item pondérer les formules~: les plus importantes respectent la \omterm{def:b1:lewis-resonance-vsepr:lewis}{règle
        de l'octet}, portent le moins de \omterm{def:b1:lewis-resonance-vsepr:formal-charge}{charges formelles} et placent les
        charges négatives sur les atomes les plus électronégatifs~; des
        formules équivalentes ont le même poids.
\end{enumerate}
\end{method}

\begin{example}[Ozone, nitrate, benzène]\label{ex:b1:lewis-resonance-vsepr:examples}
L'ozone a deux formules équivalentes~: chaque liaison \ce{O-O} est double
dans l'une et simple dans l'autre, si bien que toutes deux ont un indice
de liaison $1.5$ et la même longueur, \qty{127.8}{pm}, entre \ce{O=O} et
\ce{O-O}. L'ion nitrate \ce{NO3-} a trois formules équivalentes~; chaque
liaison \ce{N-O} a pour indice $4/3$. Le benzène, \ce{C6H6}, a deux
formules de Kekulé~; ses six liaisons \ce{C-C} sont égales,
\qty{139.7}{pm}, entre la liaison simple de l'éthane (\qty{153.6}{pm})
et la liaison double de l'éthène (\qty{133.9}{pm}).
% ledger: geo:O3.rOO, geo:C6H6.rCC, geo:C2H6.rCC, geo:C2H4.rCC
\end{example}

\begin{omfigure}
\vspace*{10pt}
\large
\setchemfig{atom sep=2.4em}
\schemestart
\chemfig{\charge{135=\|,225=\|}{O}=[:-30]\charge{270=\|,90:2pt=$\scriptstyle\oplus$}{O}-[:30]\charge{90=\|,0=\|,270=\|,45:5pt=$\scriptstyle\ominus$}{O}}
\arrow{<->}
\chemfig{\charge{90=\|,180=\|,270=\|,135:5pt=$\scriptstyle\ominus$}{O}-[:-30]\charge{270=\|,90:2pt=$\scriptstyle\oplus$}{O}=[:30]\charge{45=\|,315=\|}{O}}
\schemestop

\medskip
\schemestart
\chemfig{*6(=-=-=-)}
\arrow{<->}
\chemfig{*6(-=-=-=)}
\schemestop
\hspace{3em}
\chemfig{\charge{270:2pt=$\scriptstyle\oplus$}{N}(=[:90]\charge{45=\|,135=\|}{O})(-[:210]\charge{150=\|,240=\|,300=\|,90:2pt=$\scriptstyle\ominus$}{O})-[:330]\charge{30=\|,300=\|,240=\|,90:2pt=$\scriptstyle\ominus$}{O}}
\normalsize
\par\vspace{14pt}
{\small Mésomérie dans l'ozone (un \omterm{def:b1:lewis-resonance-vsepr:lewis}{doublet non liant} de l'oxygène de
droite devient une liaison $\pi$ tandis que la liaison $\pi$ de gauche
devient un \omterm{def:b1:lewis-resonance-vsepr:lewis}{doublet non liant}) et dans le benzène, et l'une des trois
formules équivalentes de l'ion nitrate, dont la charge est partagée par
les trois oxygènes.}
\end{omfigure}

\begin{definition}[Liaison $\sigma$, liaison $\pi$]\label{def:b1:lewis-resonance-vsepr:sigma-pi}
Une liaison simple est une \emph{liaison $\sigma$}\index{liaison sigma@liaison $\sigma$}~:
son doublet se trouve sur l'axe qui joint les deux noyaux, issu du
recouvrement axial de deux orbitales. La deuxième et la troisième
liaison d'une liaison multiple sont des \emph{liaisons $\pi$}\index{liaison pi@liaison $\pi$}~:
leurs doublets se trouvent de part et d'autre de l'axe, issus du
recouvrement latéral de deux orbitales p perpendiculaires à celui-ci.
Une liaison double est formée d'une liaison $\sigma$ et d'une liaison
$\pi$~; une liaison triple d'une liaison $\sigma$ et de deux liaisons $\pi$.
\end{definition}

\begin{omfigure}
\begin{tikzpicture}[font=\small]
  % sigma: two lobes overlapping head-on
  \fill[omDef!35, opacity=0.8] (0,0) ellipse (0.9 and 0.42);
  \fill[omProp!35, opacity=0.8] (1.1,0) ellipse (0.9 and 0.42);
  \fill (0.1,0) circle (1.6pt); \fill (1.0,0) circle (1.6pt);
  \draw[dashed, black!50] (-1.2,0) -- (2.3,0);
  \node[below] at (0.55,-1.05) {$\sigma$~: recouvrement axial};
  % pi: two p orbitals side by side
  \begin{scope}[xshift=5.2cm]
    \foreach \x/\c in {0/omDef, 1.2/omProp} {
      \fill[\c!35, opacity=0.85] (\x,0.45) ellipse (0.32 and 0.45);
      \fill[\c!15, opacity=0.85] (\x,-0.45) ellipse (0.32 and 0.45);
      \draw[\c!70!black] (\x,0.45) ellipse (0.32 and 0.45) (\x,-0.45) ellipse (0.32 and 0.45);
      \fill (\x,0) circle (1.6pt);
    }
    \draw[dashed, black!50] (-0.7,0) -- (1.9,0);
    \draw[omThm, thick, densely dotted] (0.25,0.75) -- (0.95,0.75) (0.25,-0.75) -- (0.95,-0.75);
    \node[below] at (0.6,-1.05) {$\pi$~: latéral, au-dessus et au-dessous};
  \end{scope}
\end{tikzpicture}

{\small Les deux types de recouvrement. Une liaison $\sigma$ permet la
rotation autour de son axe~; une liaison $\pi$ non, ce qui rend une
liaison double rigide (\cref{ch:b1:stereochemistry-in-depth}).}
\end{omfigure}

\section{Le modèle VSEPR}

\begin{definition}[Modèle VSEPR, nombre stérique]\label{def:b1:lewis-resonance-vsepr:vsepr}
Dans le \emph{modèle VSEPR}\index{modele VSEPR@modèle VSEPR} (\textit{valence shell electron pair
repulsion}, répulsion des doublets de la couche de valence), les doublets
qui entourent un atome central A se repoussent et adoptent la
disposition qui les éloigne le plus les uns des autres. Une molécule
s'écrit \ce{AX_nE_m}, où $n$ est le nombre d'atomes liés à A (une
liaison multiple compte une fois) et $m$ le nombre de \omterm{def:b1:lewis-resonance-vsepr:lewis}{doublets non
liants} de A. Le \emph{nombre stérique}\index{nombre sterique@nombre stérique} $n + m$ fixe la disposition des
doublets~: 2 linéaire, 3 triangulaire plane, 4 tétraédrique, 5
bipyramide trigonale, 6 octaédrique. La géométrie de la molécule est la
disposition de ses seuls atomes.
\end{definition}

\begin{proposition}[Géométries VSEPR]\label{prop:b1:lewis-resonance-vsepr:geometries}
Le \omterm{def:b1:lewis-resonance-vsepr:vsepr}{modèle VSEPR} prévoit les géométries et les angles idéaux du tableau
ci-dessous~; les angles mesurés en sont proches.
\begin{center}
\small
\begin{tabular}{llllc}
\hline
type & disposition des doublets & géométrie & exemple & angle mesuré \\
\hline
\ce{AX2} & linéaire & linéaire & \ce{CO2} & $180^\circ$ \\
\ce{AX3} & triangulaire plane & triangulaire plane & \ce{BF3} & $120^\circ$ \\
\ce{AX2E} & triangulaire plane & coudée & \ce{SO2} & $119.5^\circ$ \\
\ce{AX4} & tétraédrique & tétraédrique & \ce{CH4} & $109.5^\circ$ \\
\ce{AX3E} & tétraédrique & pyramide trigonale & \ce{NH3} & $106.7^\circ$ \\
\ce{AX2E2} & tétraédrique & coudée & \ce{H2O} & $104.5^\circ$ \\
\ce{AX5} & bipyramide trigonale & bipyramide trigonale & \ce{PCl5} & $90^\circ$, $120^\circ$ \\
\ce{AX4E} & bipyramide trigonale & en bascule & \ce{SF4} & $101.6^\circ$, $173.1^\circ$ \\
\ce{AX3E2} & bipyramide trigonale & en T & \ce{ClF3} & $87.5^\circ$ \\
\ce{AX2E3} & bipyramide trigonale & linéaire & \ce{XeF2} & $180^\circ$ \\
\ce{AX6} & octaédrique & octaédrique & \ce{SF6} & $90^\circ$ \\
\ce{AX5E} & octaédrique & pyramide à base carrée & \ce{BrF5} & --- \\
\ce{AX4E2} & octaédrique & plan carré & \ce{XeF4} & --- \\
\hline
\end{tabular}
\end{center}
% ledger: geo:CO2.aOCO, geo:BF3.aFBF, geo:SO2.aOSO, geo:CH4.aHCH,
% geo:NH3.aHNH, geo:H2O.aHOH, geo:SF6.aFSF, geo:SF4.aFSF2, geo:SF4.aFSF,
% geo:ClF3.aFClF, geo:XeF2.aFXeF
\end{proposition}
\admitted

\begin{proposition}[Les doublets non liants referment les angles]\label{prop:b1:lewis-resonance-vsepr:lone-pair-angles}
Un \omterm{def:b1:lewis-resonance-vsepr:lewis}{doublet non liant}, retenu par un seul noyau, est plus étalé qu'un
\omterm{def:b1:lewis-resonance-vsepr:lewis}{doublet liant} et repousse plus fortement~: les répulsions décroissent
dans l'ordre non liant/non liant $>$ non liant/liant $>$ liant/liant.
Les angles entre liaisons diminuent donc à mesure que l'on ajoute des
\omterm{def:b1:lewis-resonance-vsepr:lewis}{doublets non liants} ($109.5^\circ$ dans \ce{CH4}, $106.7^\circ$ dans
\ce{NH3}, $104.5^\circ$ dans \ce{H2O}), et dans une bipyramide trigonale
les \omterm{def:b1:lewis-resonance-vsepr:lewis}{doublets non liants} occupent les positions équatoriales, où ils
n'ont que deux voisins à $90^\circ$.
% ledger: geo:CH4.aHCH, geo:NH3.aHNH, geo:H2O.aHOH
\end{proposition}
\admitted

\begin{omfigure}
\setchemfig{atom sep=2.6em}
\renewcommand{\arraystretch}{1.2}
\begin{tabular}{ccccc}
\chemfig{O=C=O} &
\chemfig{B(-[:90]F)(-[:210]F)(-[:330]F)} &
\chemfig{C(-[:90]H)(-[:200]H)(<[:280]H)(<:[:335]H)} &
\chemfig{P(-[:90]Cl)(-[:270]Cl)(-[:180]Cl)(<[:320]Cl)(<:[:30]Cl)} &
\chemfig{S(-[:90]F)(-[:270]F)(-[:180]F)(-[:0]F)(<[:225]F)(<:[:45]F)} \\
\small \ce{AX2}, \ce{CO2} & \small \ce{AX3}, \ce{BF3} & \small \ce{AX4}, \ce{CH4} &
\small \ce{AX5}, \ce{PCl5} & \small \ce{AX6}, \ce{SF6} \\[16pt]
\chemfig{\charge{90=\:}{S}(=[:210]O)(=[:330]O)} &
\chemfig{\charge{90=\:}{N}(-[:200]H)(<[:280]H)(<:[:335]H)} &
\chemfig{\charge{60=\:,120=\:}{O}(-[:215]H)(-[:325]H)} &
\chemfig{\charge{0=\:}{S}(-[:90]F)(-[:270]F)(<[:210]F)(<:[:150]F)} &
\chemfig{Xe(-[:0]F)(-[:90]F)(-[:180]F)(-[:270]F)} \\
\small \ce{AX2E}, \ce{SO2} & \small \ce{AX3E}, \ce{NH3} & \small \ce{AX2E2}, \ce{H2O} &
\small \ce{AX4E}, \ce{SF4} & \small \ce{AX4E2}, \ce{XeF4} \\
\end{tabular}

\medskip
{\small Géométries prévues par la méthode VSEPR. Un coin plein pointe
vers le lecteur, un coin hachuré s'en éloigne~; les liaisons simples
sont dans le plan de la page. Les \omterm{def:b1:lewis-resonance-vsepr:lewis}{doublets non liants} de l'atome central
sont représentés par des paires de points~; \ce{XeF4} est dessiné de
face, ses deux \omterm{def:b1:lewis-resonance-vsepr:lewis}{doublets non liants} pointant vers le lecteur et à
l'opposé.}
\end{omfigure}

\begin{method}[Prévoir une géométrie]\label{met:b1:lewis-resonance-vsepr:vsepr}
Pour prévoir la géométrie autour d'un atome central~:
\begin{enumerate}
  \item dessiner le \omterm{def:b1:lewis-resonance-vsepr:lewis}{schéma de Lewis} (\cref{met:b1:lewis-resonance-vsepr:lewis})~;
  \item compter les atomes liés à l'atome central ($n$) et ses \omterm{def:b1:lewis-resonance-vsepr:lewis}{doublets
        non liants} ($m$)~; écrire \ce{AX_nE_m}~;
  \item lire la disposition sur $n + m$ et la géométrie sur les positions
        des $n$ atomes~;
  \item corriger les angles idéaux~: les \omterm{def:b1:lewis-resonance-vsepr:lewis}{doublets non liants} et les
        liaisons multiples prennent plus de place que les liaisons simples.
\end{enumerate}
\end{method}

\begin{definition}[Hybridation]\label{def:b1:lewis-resonance-vsepr:hybridisation}
L'\emph{hybridation}\index{hybridation} d'un atome est une description de
ses liaisons par des orbitales pointant dans les directions prévues par
la méthode VSEPR~: un atome de \omterm{def:b1:lewis-resonance-vsepr:vsepr}{nombre stérique} 4 est dit sp$^3$ (quatre
orbitales équivalentes à $109.5^\circ$), de \omterm{def:b1:lewis-resonance-vsepr:vsepr}{nombre stérique} 3 sp$^2$
(trois orbitales à $120^\circ$ dans un plan, une orbitale p restant
perpendiculaire à ce plan), de \omterm{def:b1:lewis-resonance-vsepr:vsepr}{nombre stérique} 2 sp (deux orbitales à
$180^\circ$, deux orbitales p restantes). Les orbitales p restantes
forment les liaisons $\pi$.
\end{definition}

\begin{example}[Les carbones de la chimie organique]\label{ex:b1:lewis-resonance-vsepr:carbon}
Dans l'éthane, chaque carbone est sp$^3$, tétraédrique~; dans l'éthène
sp$^2$, les six atomes dans un même plan avec des angles voisins de
$120^\circ$ (\ce{H-C-H} mesuré à $117.6^\circ$)~; dans l'éthyne sp, les
quatre atomes alignés. Les longueurs de liaison diminuent de 153.6 à
133.9 puis à \qty{120.3}{pm} à mesure que s'ajoutent des liaisons $\pi$.
% ledger: geo:C2H4.aHCH, geo:C2H6.rCC, geo:C2H4.rCC, geo:C2H2.rCC
\end{example}

\begin{remark}[Une description, pas une cause]\label{rem:b1:lewis-resonance-vsepr:hybrid-limit}
L'\omterm{def:b1:lewis-resonance-vsepr:hybridisation}{hybridation} décrit une géométrie déjà connue~; elle ne l'explique pas,
et elle échoue pour les énergies des électrons, dont rend compte la
théorie des orbitales moléculaires, dans le volume de deuxième année.
Elle reste le langage courant de la chimie organique~: un «~carbone
sp$^2$~» désigne un carbone plan doté d'une orbitale p libre pour une
liaison $\pi$.
\end{remark}

\section{Molécules polaires~: moments dipolaires}

\begin{definition}[Moment dipolaire de liaison, moment dipolaire]\label{def:b1:lewis-resonance-vsepr:dipole}
Dans une liaison entre atomes d'\omterm{def:b1:periodicity:electronegativity}{électronégativités} différentes, les
électrons penchent vers l'atome le plus électronégatif, qui porte une
charge partielle $-\delta e$ tandis que son partenaire porte $+\delta e$
($0 < \delta < 1$). Le \emph{moment dipolaire de liaison}\index{moment dipolaire de liaison} est le
vecteur $\vec\mu = \delta e\,
\vec d$, de norme $\delta e d$, orienté de la charge partielle négative
vers la charge partielle positive, $\vec d$ étant le vecteur qui les
joint. Le \emph{moment dipolaire}\index{moment dipolaire} d'une molécule est la somme
vectorielle de ses moments de liaison (et des contributions de ses
\omterm{def:b1:lewis-resonance-vsepr:lewis}{doublets non liants})~; il s'exprime en coulombs-mètres ou en debyes,
$\qty{1}{\debye} =
\qty{3.336e-30}{C.m}$. Une molécule de moment dipolaire non nul est une
\emph{molécule polaire}\index{molecule polaire@molécule polaire}.
% ledger: const:c (1 D = 1e-21/c C.m)
\end{definition}

\begin{remark}[Convention d'orientation]\label{rem:b1:lewis-resonance-vsepr:convention}
La physique oriente le vecteur dipolaire de la charge négative vers la
charge positive~; c'est la convention retenue ici. Des ouvrages de
chimie plus anciens dessinent une flèche à queue barrée pointant vers
l'extrémité négative~; seul le sens diffère.
\end{remark}

\begin{proposition}[Moment dipolaire d'une molécule symétrique]\label{prop:b1:lewis-resonance-vsepr:dipole-sum}
Une molécule \ce{AX_n} sans \omterm{def:b1:lewis-resonance-vsepr:lewis}{doublet non liant} sur A (\ce{AX2} linéaire,
\ce{AX3} triangulaire, \ce{AX4} tétraédrique, \dots) a un \omterm{def:b1:lewis-resonance-vsepr:dipole}{moment
dipolaire} nul, quelle que soit la polarité de ses liaisons. Une molécule
\ce{AX2} coudée d'angle $\theta$, dont les liaisons ont le moment
$\mu_b$, a pour \omterm{def:b1:lewis-resonance-vsepr:dipole}{moment dipolaire}
\[
  \mu = 2\mu_b \cos\frac{\theta}{2} .
\]
\end{proposition}

\begin{proof}
Dans les géométries symétriques, les vecteurs de liaison ont même norme
et des directions disposées symétriquement autour de A~: leur somme est
nulle (pour \ce{AX2} linéaire ils sont opposés~; pour \ce{AX3} à
$120^\circ$, trois vecteurs de même norme ont une somme nulle~; pour
\ce{AX4}, la somme est invariante par les rotations du tétraèdre, donc
nulle). Pour la molécule coudée, chaque vecteur de liaison fait l'angle
$\theta/2$ avec la bissectrice~; les composantes perpendiculaires à la
bissectrice s'annulent et celles qui lui sont parallèles s'ajoutent~:
$2\mu_b\cos(\theta/2)$.
\end{proof}

\begin{omfigure}
\begin{tikzpicture}[font=\small, >={Stealth}]
  % water
  \node[aO] (o) at (0,0) {};
  \node[aH] (h1) at (-0.95,-0.73) {};
  \node[aH] (h2) at (0.95,-0.73) {};
  \begin{scope}[on background layer]
    \draw[bond] (o.center) -- (h1.center); \draw[bond] (o.center) -- (h2.center);
  \end{scope}
  \draw[->, omProp, thick] (-0.25,0.15) -- (-0.95,-0.38);
  \draw[->, omProp, thick] (0.25,0.15) -- (0.95,-0.38);
  \draw[->, omThm, very thick] (0,-0.25) -- (0,-1.45);
  \node[omThm, right] at (0.05,-1.25) {$\vec\mu$};
  \node at (0,0.55) {$\delta^-$};
  \node at (-1.25,-1.05) {$\delta^+$}; \node at (1.25,-1.05) {$\delta^+$};
  \node[below] at (0,-1.7) {\ce{H2O}~: $\mu = \qty{1.86}{\debye}$};
  % carbon dioxide
  \begin{scope}[xshift=5cm]
    \node[aO] (a) at (-1.1,0) {}; \node[aC] (c) at (0,0) {}; \node[aO] (b) at (1.1,0) {};
    \begin{scope}[on background layer]
      \foreach \d in {-1.6,1.6} {
        \draw[bond, line width=1.1pt] ([yshift=\d pt]a.center) -- ([yshift=\d pt]c.center);
        \draw[bond, line width=1.1pt] ([yshift=\d pt]c.center) -- ([yshift=\d pt]b.center);}
    \end{scope}
    \draw[->, omProp, thick] (-0.75,0.45) -- (-0.2,0.45);
    \draw[->, omProp, thick] (0.75,0.45) -- (0.2,0.45);
    \node[below] at (0,-1.7) {\ce{CO2}~: moments compensés, $\mu = 0$};
  \end{scope}
\end{tikzpicture}

{\small \omterm{def:b1:lewis-resonance-vsepr:dipole}{Moments de liaison} (en orange, de la charge partielle négative
vers la charge partielle positive) et leur somme (en rouge). Dans l'eau,
coudée, ils s'ajoutent~; dans le dioxyde de carbone, linéaire, ils se
compensent.}
% ledger: dip:H2O
\end{omfigure}

\begin{example}[Molécules polaires et apolaires]\label{ex:b1:lewis-resonance-vsepr:dipoles}
\ce{CO2}, \ce{BF3}, \ce{CH4} et \ce{CCl4} ont un \omterm{def:b1:lewis-resonance-vsepr:dipole}{moment dipolaire} nul.
\ce{H2O} (\qty{1.86}{\debye}), \ce{NH3} (\qty{1.48}{\debye}) et
\ce{SO2} (\qty{1.63}{\debye}) sont polaires. De \ce{CH3Cl} à
\ce{CH2Cl2} puis \ce{CHCl3}, le moment diminue, 1.87, 1.62,
\qty{1.04}{\debye}, pour s'annuler dans \ce{CCl4}~: les \omterm{def:b1:lewis-resonance-vsepr:dipole}{moments des
liaisons} \ce{C-Cl} se compensent de plus en plus. Chez les halogénures
d'hydrogène, il suit la différence d'\omterm{def:b1:periodicity:electronegativity}{électronégativité}~: HF 1.83, HCl
1.09, HBr 0.83, HI \qty{0.45}{\debye}.
% ledger: dip:H2O, dip:NH3, dip:SO2, dip:CH3Cl, dip:CH2Cl2, dip:CHCl3,
% dip:CCl4, dip:HF, dip:HCl, dip:HBr, dip:HI
\end{example}

\section{Exercices}

\begin{exercise}[$\star$]\label{exo:b1:lewis-resonance-vsepr:1}
Dessiner les \omterm{def:b1:lewis-resonance-vsepr:lewis}{schémas de Lewis} de \ce{H2O}, \ce{NH3}, \ce{CO2}, \ce{HCN},
\ce{CH2O} (méthanal) et \ce{NH4+}, et donner les \omterm{def:b1:lewis-resonance-vsepr:formal-charge}{charges formelles}.
\end{exercise}

\begin{exercise}[$\star$]\label{exo:b1:lewis-resonance-vsepr:2}
Donner le type \ce{AX_nE_m}, la géométrie et l'angle approximatif de
\ce{H2S}, \ce{PCl3}, \ce{BF3}, \ce{SiH4} et \ce{CO3^{2-}} (l'atome
central figurant en premier dans chaque formule).
\end{exercise}

\begin{exercise}[$\star$]\label{exo:b1:lewis-resonance-vsepr:3}
Lesquelles de ces molécules sont polaires~: \ce{CCl4}, \ce{CH2Cl2},
\ce{BF3}, \ce{NH3}, \ce{CO2}, \ce{SO2}~? Justifier à partir des géométries.
\end{exercise}

\begin{exercise}[$\star$]\label{exo:b1:lewis-resonance-vsepr:4}
Identifier l'\omterm{def:b1:lewis-resonance-vsepr:lewis-acid}{acide de Lewis} et la \omterm{def:b1:lewis-resonance-vsepr:lewis-acid}{base de Lewis} dans \ce{H+ + H2O -> H3O+}
et dans \ce{Cu^{2+} + 4NH3 -> [Cu(NH3)4]^{2+}}. Quelles liaisons sont
datives~?
\end{exercise}

\begin{exercise}[$\star\star$]\label{exo:b1:lewis-resonance-vsepr:5}
Écrire deux \omterm{def:b1:lewis-resonance-vsepr:resonance}{formules mésomères} de l'ion éthanoate \ce{CH3COO-}. Que
prévoient-elles pour les longueurs des deux liaisons \ce{C-O}~? Comparer
à l'acide méthanoïque \ce{HCOOH}, dont les deux liaisons \ce{C-O}
mesurent 120.2 et \qty{134.3}{pm}.
% ledger: geo:H2CO2.rCO, geo:H2CO2.rCO2
\end{exercise}

\begin{exercise}[$\star\star$]\label{exo:b1:lewis-resonance-vsepr:6}
Dessiner le \omterm{def:b1:lewis-resonance-vsepr:lewis}{schéma de Lewis} de l'ion sulfate \ce{SO4^{2-}} avec quatre
liaisons simples \ce{S-O}, calculer les \omterm{def:b1:lewis-resonance-vsepr:formal-charge}{charges formelles} et prévoir la
géométrie. Combien de \omterm{def:b1:lewis-resonance-vsepr:resonance}{formules mésomères} comportant deux liaisons
\ce{S=O} et aucune charge sur le soufre peut-on écrire~?
\end{exercise}

\begin{exercise}[$\star\star$]\label{exo:b1:lewis-resonance-vsepr:7}
Prévoir par la méthode VSEPR les géométries de \ce{SF4}, \ce{ClF3},
\ce{XeF2} et \ce{XeF4}. Expliquer pourquoi les \omterm{def:b1:lewis-resonance-vsepr:lewis}{doublets non liants} d'une
bipyramide trigonale se placent dans le plan équatorial.
\end{exercise}

\begin{exercise}[$\star\star$]\label{exo:b1:lewis-resonance-vsepr:8}
Les angles \ce{H-X-H} valent $104.5^\circ$ dans \ce{H2O} et $92.1^\circ$
dans \ce{H2S}~; $106.7^\circ$ dans \ce{NH3} et $93.3^\circ$ dans
\ce{PH3}. Proposer une explication fondée sur la taille et
l'\omterm{def:b1:periodicity:electronegativity}{électronégativité} de l'atome central.
% ledger: geo:H2O.aHOH, geo:H2S.aHSH, geo:NH3.aHNH, geo:PH3.aHPH
\end{exercise}

\begin{exercise}[$\star\star$]\label{exo:b1:lewis-resonance-vsepr:9}
Donner l'\omterm{def:b1:lewis-resonance-vsepr:hybridisation}{hybridation} de chaque carbone et de l'azote dans \ce{CH3-C#N}
(éthanenitrile) et dans \ce{CH2=CH-CH3}, ainsi que le nombre de liaisons
$\sigma$ et $\pi$ de chaque molécule.
\end{exercise}

\begin{exercise}[$\star\star\star$]\label{exo:b1:lewis-resonance-vsepr:10}
Le \omterm{def:b1:lewis-resonance-vsepr:dipole}{moment dipolaire} de \ce{NF3} ne vaut que \qty{0.24}{\debye}, contre
\qty{1.48}{\debye} pour \ce{NH3}, bien que la liaison \ce{N-F} soit plus
polarisée que la liaison \ce{N-H}. Expliquer la différence à l'aide des
\omterm{def:b1:lewis-resonance-vsepr:dipole}{moments de liaison} et du \omterm{def:b1:lewis-resonance-vsepr:lewis}{doublet non liant} de l'azote.
% ledger: dip:NF3, dip:NH3
\end{exercise}

\begin{exercise}[$\star\star\star$]\label{exo:b1:lewis-resonance-vsepr:11}
La molécule \ce{SO2} a un angle de $119.5^\circ$ et un \omterm{def:b1:lewis-resonance-vsepr:dipole}{moment dipolaire}
de \qty{1.63}{\debye}. Calculer le moment d'une liaison \ce{S-O}~; avec
la longueur de liaison \qty{143.2}{pm}, en déduire la charge partielle
$\delta$ portée par chaque oxygène.
% ledger: geo:SO2.aOSO, geo:SO2.rSO, dip:SO2
\end{exercise}

\begin{exercise}[$\star\star\star$]\label{exo:b1:lewis-resonance-vsepr:12}
Dessiner les \omterm{def:b1:lewis-resonance-vsepr:lewis}{schémas de Lewis} les plus importants du monoxyde de diazote
\ce{N2O} (linéaire, \ce{N-N-O}) et calculer les \omterm{def:b1:lewis-resonance-vsepr:formal-charge}{charges formelles}. Les
longueurs de liaison mesurées sont \ce{N-N} \qty{112.8}{pm} et \ce{N-O}
\qty{118.4}{pm}~; dans \ce{N2}, la liaison triple mesure
\qty{109.8}{pm}. Quelle formule a le plus de poids~?
% ledger: geo:N2O.rNN, geo:N2O.rNO, re:N2
\end{exercise}

\section{Problème~: les molécules d'un coup de foudre}

\begin{problem}[{Problème du week-end --- les espèces azotées et
oxygénées formées quand la foudre chauffe l'air~: schémas de Lewis,
mésomérie, géométries, et les charges partielles de l'eau}]
\label{pb:b1:lewis-resonance-vsepr:1}
Un éclair porte l'air à des milliers de degrés~: \ce{N2} et \ce{O2}
donnent \ce{NO}, puis \ce{NO2}, \ce{N2O}, l'ozone et finalement l'acide
nitrique des pluies. Données~: $\qty{1}{\debye} = \qty{3.336e-30}{C.m}$,
$e =
\qty{1.602e-19}{C}$. Les valeurs mesurées sont données là où il le faut.

\medskip
\noindent\textbf{Partie I --- \omterm{def:b1:lewis-resonance-vsepr:lewis}{Schémas de Lewis}.}
\begin{enumerate}
  \item Compter les \omterm{def:b1:quantum-numbers:configuration}{électrons de valence} de \ce{N2}, \ce{NO}, \ce{NO2},
        \ce{N2O}, \ce{O3} et \ce{HNO3}.
  \item Dessiner le \omterm{def:b1:lewis-resonance-vsepr:lewis}{schéma de Lewis} de \ce{N2}.
  \item Dessiner un \omterm{def:b1:lewis-resonance-vsepr:lewis}{schéma de Lewis} de \ce{NO}. Pourquoi la \omterm{def:b1:lewis-resonance-vsepr:lewis}{règle de
        l'octet} ne peut-elle être satisfaite sur les deux atomes~?
  \item Dessiner deux \omterm{def:b1:lewis-resonance-vsepr:lewis}{schémas de Lewis} de \ce{N2O}, avec \ce{N=N=O} et
        avec \ce{N#N-O}, et donner les \omterm{def:b1:lewis-resonance-vsepr:formal-charge}{charges formelles}.
  \item Dessiner un \omterm{def:b1:lewis-resonance-vsepr:lewis}{schéma de Lewis} de l'ozone et donner les \omterm{def:b1:lewis-resonance-vsepr:formal-charge}{charges
        formelles}.
  \item Dessiner un \omterm{def:b1:lewis-resonance-vsepr:lewis}{schéma de Lewis} de \ce{HNO3} et donner les \omterm{def:b1:lewis-resonance-vsepr:formal-charge}{charges
        formelles}.
  \item Lesquelles des six espèces ont un \omterm{def:b1:quantum-numbers:unpaired}{électron célibataire}~?
\end{enumerate}

\medskip
\noindent\textbf{Partie II --- Mésomérie et longueurs de liaison.}
\begin{enumerate}[resume]
  \item Écrire les deux \omterm{def:b1:lewis-resonance-vsepr:resonance}{formules mésomères} de l'ozone. Quel indice de
        liaison prévoient-elles~?
  \item La liaison \ce{O-O} mesure \qty{127.8}{pm} dans l'ozone,
        \qty{120.8}{pm} dans \ce{O2} et \qty{147.5}{pm} dans \ce{H2O2}.
        La prévision est-elle confirmée~?
  \item Écrire les trois \omterm{def:b1:lewis-resonance-vsepr:resonance}{formules mésomères} de \ce{NO3-} et donner
        l'indice de chaque liaison \ce{N-O}.
  \item Dans \ce{HNO3}, les trois liaisons \ce{N-O} mesurent 140.6, 121.1
        et \qty{119.9}{pm}. Les attribuer et expliquer à l'aide de la
        mésomérie.
  \item Expliquer pourquoi l'angle \ce{O-N-O} de \ce{NO2} ($134^\circ$)
        est supérieur à $120^\circ$.
  \item Laquelle des deux formules de la question 4 place la \omterm{def:b1:lewis-resonance-vsepr:formal-charge}{charge
        formelle} négative sur l'atome le plus électronégatif~?
\end{enumerate}

\medskip
\noindent\textbf{Partie III --- Géométries et polarité.}
\begin{enumerate}[resume]
  \item Donner le type \ce{AX_nE_m} de l'atome central de \ce{O3},
        \ce{N2O}, \ce{NO3-} et de l'azote de \ce{HNO3}, ainsi que leurs
        géométries.
  \item L'angle de l'ozone vaut $116.8^\circ$. Expliquer pourquoi il est
        inférieur à $120^\circ$.
  \item L'ozone a un \omterm{def:b1:lewis-resonance-vsepr:dipole}{moment dipolaire} de \qty{0.53}{\debye}, \ce{N2O}
        \qty{0.16}{\debye}, \ce{N2} aucun. Expliquer.
  \item Pourquoi \ce{CO2} est-il apolaire alors que \ce{SO2} est polaire~?
  \item Prévoir l'angle de \ce{NH3} par rapport à $109.5^\circ$, puis
        comparer à la valeur mesurée, $106.7^\circ$.
  \item Même question pour l'eau ($104.5^\circ$).
\end{enumerate}
% ledger: geo:O3.rOO, re:O2, geo:H2O2.rOO, geo:HNO3.rNO, geo:HNO3.rNO2,
% geo:HNO3.rNO3, geo:NO2.aONO, geo:O3.aOOO, dip:O3, dip:N2O

\medskip
\noindent\textbf{Partie IV --- Les charges partielles de l'eau.}
Le \omterm{def:b1:lewis-resonance-vsepr:dipole}{moment dipolaire} de l'eau vaut \qty{1.857}{\debye}, son angle
$104.48^\circ$ et la longueur de sa liaison \ce{O-H} \qty{95.8}{pm}.
% ledger: dip:H2O, geo:H2O.aHOH, geo:H2O.rOH
\begin{enumerate}[resume]
  \item Exprimer le \omterm{def:b1:lewis-resonance-vsepr:dipole}{moment dipolaire} de l'eau en fonction du \omterm{def:b1:lewis-resonance-vsepr:dipole}{moment de
        liaison} $\mu_{OH}$ et de l'angle.
  \item Calculer $\mu_{OH}$ en debyes.
  \item Le convertir en coulombs-mètres.
  \item Si le \omterm{def:b1:lewis-resonance-vsepr:dipole}{moment de liaison} était dû à deux charges entières $\pm e$
        placées sur les deux noyaux, que vaudrait-il~?
  \item En déduire le caractère ionique partiel $\delta$ de la liaison
        \ce{O-H}, rapport du \omterm{def:b1:lewis-resonance-vsepr:dipole}{moment de liaison} mesuré à celui de charges
        entières.
\end{enumerate}
\end{problem}
